\cotefile{64}
\scanpage{40}
\mspage{\unsure{I c}}
% archive foliation: 39

\margin{\uline{NB} $X|\underline J\simeq\Delta$ est la configuration octaédrale\ldots}

\diagram{\begin{tikzcd}[ampersand replacement=\&,column sep=large,row sep=large]
\begin{gathered}(X,\underline I)\quad
\text{\deleted{\unsure{étale}}}\ \inserted{\underline I}\ \text{rev. étale fini de degré }4\end{gathered}
\arrow[d,leftrightarrow] \& \\
\begin{gathered}
\deleted{(\underline{\unsure{J}}(\Sigma_{\underline J}),\ t\in\Gamma(\Sigma^*_{\underline J}/S),\ \underline I\ \text{\deleted{\unsure{rev.}} torseur sous})}\\
(\underline I,\ \text{d'où }\underline J,\ t\in\Gamma(\Sigma^*_{\underline J}/S))
\end{gathered}\arrow[d,leftrightarrow]\arrow[r,leftrightarrow]
\&\begin{gathered}
\text{\uline{NB}}\quad\underline t\text{ et }\underline I\text{ indépendants}\\
(\underline J\ (\text{d'où }\underline V_{\underline J}\text{ et }\Sigma_J),\
\text{section }\inserted{t}\text{ de }\Sigma^*_J\text{ sur }S,\\
\text{torseur }I\text{ sur }\underline V_{\underline J})
\end{gathered}\\
\begin{gathered}
(\underline J,\ t\in\Gamma(\Sigma^*_{\underline J}/S),\
X\text{ rev. de }\Sigma^{\unsure{*}}_J,\ \deleted{\unsure{p}}\
\text{étale sur }\Sigma^*_J,\ \text{propre lisse sur }S,\\
\text{en chaque fibre géom. rev. \uline{abélien}}
\ \deleted{\text{\unsure{de degré } }4})\\
\inserted{\text{\unsure{fin. ram.}, \unsure{ramifiés, d’ordre} }2}\\
\text{\uline{NB}}\quad X\text{ et }t\text{ sont indép. l'un de l'autre}
\end{gathered}\arrow[d,leftrightarrow]\&\\
\begin{gathered}
(\underline J,\ t\in\Gamma(\Sigma^*_{\underline J}/S),\
\deleted{\unsure{I}}\underline K\text{ torseur sous }\underline J)\\
\text{\uline{NB}}\quad\underline K\text{ et }t\text{ sont indép. l'un de l'autre.}
\end{gathered}\&
\end{tikzcd}}

\margin{\deleted{$\unsure{\varphi_i\ \varphi_j\ \varphi_k}$}\quad
$\unsure{x_i}\quad x_j\quad x_k$
\diagram{\begin{tikzcd}[ampersand replacement=\&]
V_{\underline J}(M_0)\arrow[r]\&\unsure{\mathbb M_2}(M_0)\arrow[r,phantom,"\simeq"]\&\underline{\operatorname{End}}(M_0)\\
\&\&1\arrow[u,phantom,"\in" description]
\end{tikzcd}}}

\medskip\hrule\medskip
Quand $\underline I$ est muni d'une section, de sorte que $\underline I\simeq\underline J\amalg S$, alors on a un iso. can. $X\simeq\Sigma_J$, d'où $f_{\deleted{\unsure{I}}}:\Sigma_J\mathbin{\unsure{-|\longrightarrow}}\Sigma_J$.
